\documentclass[11pt,a4paper]{article}
\usepackage[margin=27mm]{geometry}
\usepackage{amsmath,amssymb,amsthm,fontspec,unicode-math,hyperref,booktabs,longtable}
\usepackage{array,calc,fvextra}
\setmainfont{STIXGeneral}
\setmathfont{STIX Two Math}
\setmonofont{Menlo}
\hypersetup{hidelinks}
\setlength{\emergencystretch}{2em}
\providecommand{\tightlist}{\setlength{\itemsep}{0pt}\setlength{\parskip}{0pt}}
\newtheorem{theorem}{Theorem}
\newtheorem{lemma}[theorem]{Lemma}
\newtheorem{corollary}[theorem]{Corollary}
\newcommand{\gen}[1]{\langle #1\rangle}
\newcommand{\Sym}{\operatorname{Sym}}
\newcommand{\Aut}{\operatorname{Aut}}
\newcommand{\Fix}{\operatorname{Fix}}

\title{Odd kernels, cyclotomic modules and cyclic carries}
\author{The human\_compression project}\date{Topic F: expanded common theory, private draft}
\begin{document}\maketitle
\begin{abstract}
We combine the complete existing odd-kernel and cyclotomic lattice
arguments with a cyclic-extension mechanism whose nonzero carry
obstructs reversal. A short lower-central calculation handles all
odd primes in the small-order problem, and an invariant cyclic
kernel transfers chirality in the previously studied lift. The
regular and chiral applications share the full module calculations;
each retains its own exact source scope and classical inputs.
\end{abstract}
\section{Marked groups and intersections}
% BEGIN UF.Conventions
We use $x^g=g^{-1}xg$ and $[x,y]=x^{-1}y^{-1}xy$.
A rank-three string marking of a finite group $G$ is a generating triple
of distinct nonidentity involutions $r_0,r_1,r_2$ whose endpoints commute.
Write $G_I=\gen{r_i:i\in I}$ for $I\subseteq\{0,1,2\}$.
The marking is a \emph{string C-group} if
\[
 G_I\cap G_J=G_{I\cap J}\qquad(I,J\subseteq\{0,1,2\}).
\]
Its type is $\{k_1,k_2\}$, where $k_1=|r_0r_1|$ and
$k_2=|r_1r_2|$. We call it nondegenerate when $k_1,k_2\geq3$,
and tight when $|G|=2k_1k_2$. A marked isomorphism preserves the
three generators; duality reverses them. We retain the classical
correspondence between string C-groups and abstract regular polyhedra
when using geometric terminology.

% END UF.Conventions

% BEGIN UF.RankThree
\begin{lemma}[Rank-three intersection test]\label{UF.RankThree}
A nondegenerate rank-three string marking is a string C-group if and
only if
\begin{equation}\label{eq:rankthree}
 \gen{r_0,r_1}\cap\gen{r_1,r_2}=\gen{r_1}.
\end{equation}
In that case $2k_1k_2\leq |G|$.
\end{lemma}
\begin{proof}
Necessity is part of the intersection property. For sufficiency set
$H_{ij}=\gen{r_i,r_j}$. The endpoint group consists of
$1,r_0,r_2,r_0r_2$. If its intersection with $H_{01}$ contains $r_2$
or $r_0r_2$, then $r_2\in H_{01}\cap H_{12}=\gen{r_1}$,
contrary to distinctness. Hence
$H_{02}\cap H_{01}=\gen{r_0}$; dually
$H_{02}\cap H_{12}=\gen{r_2}$. These equalities and
\eqref{eq:rankthree} also exclude a generator from the group on the
other two. Singleton intersections follow, and pairs involving the
empty or full index set are immediate. This accounts for every pair
of subsets. Finally $|H_{01}|=2k_1$ and $|H_{12}|=2k_2$, so
\[
 |H_{01}H_{12}|=\frac{|H_{01}|\,|H_{12}|}{|H_{01}\cap H_{12}|}
 =2k_1k_2\leq |G|.
\]
\end{proof}
% END UF.RankThree

% BEGIN UF.PrimeIndex
\begin{lemma}[Prime-index intersection]\label{UF.PrimeIndex}
Let $C,H,J$ be subgroups of a group, with $C\leq H\cap J$ and
$[H:C]=p$ prime. If $H\not\leq J$, then $H\cap J=C$.
In particular, this holds if $\gen{H,J}=G$ and $J<G$.
\end{lemma}
\begin{proof}
The index tower gives
$p=[H:H\cap J][H\cap J:C]$. The first factor is not one,
so the second is one. In the final situation $H\leq J$ would give
$G=\gen{H,J}=J$, contrary to the hypothesis.
\end{proof}
For later products it is useful to record the elementary identities
\begin{equation}\label{eq:reflections}
 (sh)^2=1,\qquad (sh)(bh)=sb^{-1},
\end{equation}
whenever $h^2=1$, $hsh=s^{-1}$ and $hbh=b^{-1}$.
The first equality uses only the condition involving $s$; the second
uses only the condition involving $b$.
% END UF.PrimeIndex

\section{An inverted orbit and its characters}
% BEGIN UF.Orbit
\begin{lemma}[One inverted orbit]\label{UF.Orbit}
Suppose $G=P\rtimes K$, where $P$ has odd order and
$K=\gen{a,b}$ is a Klein four group. Suppose $G=\gen{K,r}$,
where $r$ is a nonidentity involution. Write $r=ut$ with
$u\in P$, $t\in K$. Then $t\ne1$, $u^t=u^{-1}$, and
\[
 P=\gen{u,u^k}\qquad\text{for every }k\in K\setminus\gen{t}.
\]
\end{lemma}
\begin{proof}
An odd-order group has no nonidentity involution, so $t\ne1$.
The equation $(ut)^2=1$ gives $u^t=u^{-1}$. Put
$L=\gen{u,u^k}$. Since $t$ commutes with $k$, it inverts both
displayed generators of $L$. The element $k$ exchanges them. Thus
$L$ is $K$-invariant and $LK$ is a subgroup containing $K$ and $r$.
It equals $G$, and intersecting with $P$ gives $L=P$.
\end{proof}
This proof says that $t$ inverts the two indicated generators. It does
not assert that inversion is an automorphism of a nonabelian group.
% END UF.Orbit

% BEGIN UF.Characters
\begin{lemma}[Two cyclic characters]\label{UF.Characters}
In Lemma~\ref{UF.Orbit}, suppose $P$ is abelian and noncyclic.
Then it is the direct product of two nontrivial cyclic subgroups
$P_\chi$, indexed by the two characters
$\chi:K\longrightarrow\{1,-1\}$ with $\chi(t)=-1$.
The projection of $u$ generates each $P_\chi$.
\end{lemma}
\begin{proof}
Write $P$ additively. Multiplication by four is an automorphism of $P$.
The four endomorphisms
\[
 e_{\epsilon,\delta}=\tfrac14(1+\epsilon a)(1+\delta b),
 \qquad\epsilon,\delta\in\{1,-1\},
\]
are orthogonal idempotents with sum the identity: expand the products
using $a^2=b^2=1$ and $ab=ba$. Their images give the four character
components, on which $a,b$ act by $\epsilon,\delta$ respectively.
The element $t$ inverts $u$ and commutes with $K$, so it inverts the
whole $K$-orbit of $u$. That orbit generates $P$ by
Lemma~\ref{UF.Orbit}; hence $t$ acts as $-1$ on $P$.
Every component with $\chi(t)=1$ is therefore annihilated by two and
is zero. For either remaining character, the projection of each orbit
element is plus or minus $e_\chi u$. Its component is consequently
cyclic. If one component were zero then $P$ would be cyclic, so both
are nontrivial.
\end{proof}
% END UF.Characters

\section{The complete cyclotomic lattice}
Let $p$ be an odd prime, let $\zeta$ be a primitive $p$th root of unity,
and put $O=\mathbb Z[\zeta]$ and $\pi=\zeta-1$.
% BEGIN UF.Lattice
\begin{theorem}[Cyclotomic coordinates]\label{UF.Lattice}
For every $n\geq1$, the ring $A_n=O/\pi^nO$ has $p^n$ elements.
Write uniquely
\[
 n=(e-1)(p-1)+r,\qquad e\geq1,\quad 1\leq r\leq p-1.
\]
As an additive group with the indicated ordered generators,
\begin{equation}\label{eq:coordinates}
 A_n=\gen{1}\oplus\gen{\pi}\oplus\cdots\oplus\gen{\pi^{p-2}}
 \cong (\mathbb Z/p^e\mathbb Z)^r
 \oplus(\mathbb Z/p^{e-1}\mathbb Z)^{p-1-r}.
\end{equation}
A modulus-one factor is trivial. Complex conjugation induces an
involutory automorphism of $A_n$. Multiplication by $\zeta$ has order
$p$ when $n\geq2$, and for $j\not\equiv0\pmod p$ it satisfies
$\sum_{i=0}^{p-1}\zeta^{ij}=0$.
\end{theorem}
\begin{proof}
The polynomial
\[
 \Phi_p(1+X)=p+\binom p2X+\cdots+
 \binom p{p-1}X^{p-2}+X^{p-1}
\]
is Eisenstein at $p$. Thus $O$ is an integral domain with
$\mathbb Z$-basis $1,\pi,\ldots,\pi^{p-2}$. Substitution
$\zeta=1$ gives $O/\pi O\cong\mathbb F_p$.
Multiplication by $\pi^j$ induces an isomorphism
$O/\pi O\longrightarrow\pi^jO/\pi^{j+1}O$:
surjectivity is immediate and injectivity follows by cancellation
in the domain. The successive indices are all $p$, whence
\begin{equation}\label{eq:latticeorder}
 |O/\pi^nO|=p^n.
\end{equation}

The identity $p=\prod_{j=1}^{p-1}(1-\zeta^j)$ gives
$pO=\pi^{p-1}O$. Here equality, rather than inclusion alone, follows
because each ratio $(1-\zeta^j)/(1-\zeta)$ is a unit of $O$.
It belongs to $O$ by a geometric sum. The automorphism
$\zeta\mapsto\zeta^j$ identifies the norms of numerator and
denominator, both equal to $p$, so the ratio has norm one.
Its multiplication matrix on the integral basis has determinant one;
the adjugate gives an integral inverse and hence a unit of $O$.

For $1\leq r\leq p-1$ define the integral lattice
\[
 L_r=p\mathbb Z\,1+\cdots+p\mathbb Z\,\pi^{r-1}
       +\mathbb Z\,\pi^r+\cdots+\mathbb Z\,\pi^{p-2}.
\]
Every displayed generator is in $\pi^rO$, using
$pO=\pi^{p-1}O$ for the first group of generators.
Both $L_r$ and $\pi^rO$ have index $p^r$ in $O$, the latter by
\eqref{eq:latticeorder}; thus $L_r=\pi^rO$.
Now $\pi^nO=p^{e-1}\pi^rO$, which gives every coordinate modulus
in \eqref{eq:coordinates}. In particular it gives the individual
orders, including trivial listed generators when $e=1$.

The equation $\overline\pi=-\zeta^{-1}\pi$ shows that conjugation
preserves $\pi^nO$. Multiplication by $\zeta$ has $p$th power the
identity. For $n\geq2$ it is not the identity, since its difference
from the identity sends $1$ to $\pi\notin\pi^nO$; the latter follows
from the strict successive indices above. Its order is therefore $p$.
Finally $1+\zeta+\cdots+\zeta^{p-1}=0$; multiplication of exponents
by a nonzero residue $j$ permutes these terms.
\end{proof}
% END UF.Lattice

% BEGIN UF.Conjugation
\begin{corollary}[Conjugation and displacement]\label{UF.Conjugation}
Let $A_n$ be written additively and let an element $\beta$ of order
$p$ act by $v^\beta=\zeta v$. Put $s_i=\pi^{i-1}$ for $i\geq1$.
Then $[s_i,\beta]=s_{i+1}$ in $A_n\rtimes\gen{\beta}$, and
\begin{align}
 \overline{s_i}&=(-1)^{i-1}\zeta^{1-i}s_i,
       &&i\geq1,\label{eq:conjugation}\\
 \zeta^{1-i}s_i&=\sum_{j=i}^{p+1}
       \binom{p+1-i}{j-i}s_j,
       &&1\leq i\leq p.\label{eq:binomial}
\end{align}
The first expression on the left of the commutator identity is group
notation; the two displayed equations are additive ring identities.
\end{corollary}
\begin{proof}
The commutator is $(\zeta-1)s_i=\pi s_i$.
Equation~\eqref{eq:conjugation} follows by raising
$\overline\pi=-\zeta^{-1}\pi$ to the power $i-1$.
For \eqref{eq:binomial}, use $\zeta^{1-i}=\zeta^{p+1-i}$ and
expand $(1+\pi)^{p+1-i}\pi^{i-1}$. No independence of the
additional $s_p,s_{p+1}$ is asserted or needed.
\end{proof}
% END UF.Conjugation

\section{Rotation pairs and reversal}
Suppose $G=\gen{a,b}$, with $|a|,|b|\geq3$,
$(ab)^2=1$ and $\gen a\cap\gen b=1$.
If an involutory automorphism $\delta$ inverts both generators,
form $R=G\rtimes\gen\delta$ and put
\[
r_0=a\delta,\qquad r_1=\delta,\qquad r_2=\delta b.
\]
These are three distinct nonidentity involutions, the endpoints
commute, and they generate $R$. The adjacent subgroups are
$\gen a\rtimes\gen\delta$ and $\gen b\rtimes\gen\delta$;
their intersection is exactly $\gen\delta$ by the disjoint
cyclic subgroups in $G$. The complete rank-three intersection test
of Topic F proves the string C-group property. Thus $R$ is the
full regular group with the asserted type and order $2|G|$.
This gives the complete group-theoretic regular-doubling interface;
the classical string C-group--polytope correspondence is retained.


\begin{lemma}[Cyclic quotient criterion]\label{UF.CyclicQuotient}
Let $\phi:G\to H$ be a homomorphism and let $a,b\in G$.
If $\gen{\phi(a)}\cap\gen{\phi(b)}=1$ and the restriction of
$\phi$ to either $\gen a$ or $\gen b$ is injective, then
$\gen a\cap\gen b=1$.
\end{lemma}
\begin{proof}
Any common element maps into the trivial intersection downstairs.
It is therefore the identity by injectivity on the specified cyclic
subgroup. This is the quotient criterion used in source Proposition 2.1;
no finite-order or surjectivity hypothesis is needed for this step.
\end{proof}

\begin{lemma}[Abelian kernel reversal]\label{lem:abelian}
Suppose $G=P\rtimes\gen t$, where $P$ is abelian and $t^2=1$.
If $a\in P$ and $b=a^{-1}t$, an involutory automorphism of $G$
simultaneously inverts $a,b$.
\end{lemma}
\begin{proof}
Inversion on $P$ is an automorphism and commutes with the action
of $t$, so extends to an automorphism $\iota$ fixing $t$.
Follow it by conjugation with $a$. The result sends $a$ to $a^{-1}$
and $t$ to $a^{-1}ta$, hence sends $b=a^{-1}t$ to
$a(a^{-1}ta)=ta=b^{-1}$. Its square fixes $a$ and every element
of $P$, and fixes $t$ because $\iota(a)=a^{-1}$; thus it is
involutory. No presentation-order argument is needed.
\end{proof}


\section{A shared split cyclotomic construction}
% BEGIN UF.SplitConstruction
\begin{theorem}[The cyclotomic polyhedra]\label{UF.SplitConstruction}
For every odd prime $p$ and every $m\geq3$, there is a string
C-group of order $4p^m$ and type $\{p,2p\}$ with nonabelian
Sylow $p$-subgroup. More precisely, let
\[
 O=\mathbb Z[\zeta_p],\quad \pi=\zeta_p-1,\quad
 A=O/\pi^{m-1}O,\quad P=A\rtimes\gen{\beta},
 \qquad |\beta|=p,\quad v^\beta=\zeta_pv.
\]
Define automorphisms by
\[
 \sigma(v)=\overline v,\quad\sigma(\beta)=\beta^{-1},
 \qquad\tau(v)=-v,\quad\tau(\beta)=\beta.
\]
Then $G=P\rtimes\gen{\sigma,\tau}$ has the required marking
\[
 r_0=s\tau\sigma,\qquad r_1=\beta\tau\sigma,\qquad r_2=\sigma,
\]
where $s$ denotes the additive class of the ring element $1$.
It is not the identity of the additive group $A$.
\end{theorem}
\begin{proof}
Set $n=m-1\geq2$ and write $\zeta=\zeta_p$. To fix the action
convention, use pairs $(v,k)\in A\times\mathbb Z/p\mathbb Z$ with
\begin{equation}\label{eq:pairlaw}
 (v,k)(w,l)=(v+\zeta^{-k}w,k+l).
\end{equation}
Then $\beta=(0,1)$ has $v^\beta=\zeta v$.
The cyclotomic coordinates of Theorem~\ref{UF.Lattice} give $|P|=p^{n+1}=p^m$.
The action by $\zeta$ is nontrivial for $n\geq2$, so $P$ is
nonabelian. Every element outside $A$ has order exactly $p$:
for $k\ne0$ its $p$th power has first coordinate
$\sum_{i=0}^{p-1}\zeta^{-ik}v=0$, and its image in $P/A$ has
order $p$.

On pairs the proposed maps are
$\sigma(v,k)=(\overline v,-k)$ and $\tau(v,k)=(-v,k)$.
They preserve~\eqref{eq:pairlaw}, since conjugation sends
$\zeta^{-k}$ to $\zeta^k$. They are commuting involutory
automorphisms. Moreover, $\sigma$ inverts the nontrivial element
$\beta$, whereas $\tau$ fixes it; and $\tau$ is nontrivial on
the nonzero odd-order group $A$. Thus they are independent and
generate a Klein four group. In particular $|G|=4p^m$.

Put $h=\tau\sigma$. In multiplicative group notation
$s^h=s^{-1}$ and $\beta^h=\beta^{-1}$. The elementary reflection
identities above make $r_0=sh$, $r_1=\beta h$ and
$r_2=\sigma$ involutions. Directly,
\[
 r_0r_2=s\tau=r_2r_0,\qquad
 r_0r_1=s\beta^{-1},\qquad r_1r_2=\beta\tau.
\]
The middle product is outside $A$ and has order $p$.
The last has order $2p$, since its two factors commute and have
coprime orders. The marked generators are nonidentity and distinct;
the first two project to $\tau\sigma$ in the complement, the third
to $\sigma$, and the first two are distinct because $s\ne\beta$.

They generate the whole group. Indeed,
$(r_1r_2)^p=\tau$ yields $\tau$, then $\beta$; the generator
$r_2$ gives $\sigma$, and $r_0$ gives $s$. The
$\beta$-conjugates of $s$ are $1,\zeta,\ldots,\zeta^{p-1}$,
whose classes additively generate $A$.

Finally $H=\gen{r_0,r_1}$ and $J=\gen{r_1,r_2}$ have
orders $2p$ and $4p$. The subgroup $C=\gen{r_1}$ has prime
index $p$ in $H$. If $H\leq J$, all three generators would lie
in $J$, giving $G=J$ and contradicting $4p^m>4p$.
The prime-index lemma above therefore gives $H\cap J=C$.
Its rank-three test proves the full intersection property.
\end{proof}
% END UF.SplitConstruction

\section{A nonzero cyclic carry}
\begin{theorem}[A nonzero carry obstructs reversal]\label{thm:carry}
Let $p$ be an odd prime and let $M$ be a finite abelian $p$-group.
Suppose $T,S\in\operatorname{Aut}(M)$ and $s,z\in M$ satisfy
\[
T^p=S^2=1,\quad STS=T^{-1},\quad Ss=s,\quad Tz=z,\quad Sz=-z,
\quad |z|=p,\quad M=\gen{T^i s:i\in\mathbb Z}.
\]
Put $w=z+\sum_{i=0}^{p-1}T^is$, and suppose $|w|=p$ and
$\gen w\cap\gen s=0$. There is a group $P$ of order $p|M|$
containing $M$, generated by $M$ and $\beta$, with
\[
v^\beta=Tv\ (v\in M),\qquad \beta^p=z.
\]
The assignment $v\mapsto Sv$, $\beta\mapsto\beta^{-1}$ is an
involutory automorphism $\sigma$ of $P$. In $G=P\rtimes\gen\sigma$,
\[
a=\beta s,\qquad b=s^{-1}\sigma
\]
generate $G$, have orders $p^2,2|s|$, and satisfy
$(ab)^2=1$ and $\gen a\cap\gen b=1$. No automorphism of $G$
simultaneously inverts $a,b$.
\end{theorem}
\begin{proof}
On $M\times\mathbb Z$ use multiplication
$(u,i)(v,j)=(u+T^{-i}v,i+j)$. The element $(-z,p)$ is central:
$T^p=1$ and $Tz=z$. Quotient by its infinite cyclic subgroup.
The map $M\to P$ is injective, since an element $(-kz,kp)$ has
second coordinate zero only when $k=0$. Every coset has a unique
representative $(u,i)$ with $0\leq i<p$. This proves the order,
the action and the carry relation. The map $(u,i)\mapsto(Su,-i)$
is an involutory automorphism before taking the quotient, because
$ST^{-i}=T^iS$. It takes $(-z,p)$ to its inverse, so descends as stated.
It is nonidentity: its induced action on $P/M\cong C_p$ is inversion.

Since $b^2=-2s$, oddness gives $s\in\gen b$. We then recover
$\beta=as^{-1}$ and $\sigma=sb$. Their conjugates recover the
whole $T$-orbit of $s$, hence $G$. The product $ab=\beta\sigma$
is a nonidentity involution, and $|b|=2|s|$ because $s$ is fixed
by $\sigma$. The cyclic-extension multiplication gives $a^p=w$;
thus $|a|=p^2$. Any common element of $\gen a$ and $\gen b$ lies
in $P$, hence in $\gen{b^2}=\gen s$. Its image in $P/M$ is zero,
so it lies in $\gen{a^p}=\gen w$. The assumed intersection is zero.

Suppose $\alpha\in\operatorname{Aut}(G)$ inverted $a,b$.
The unique Sylow $p$-subgroup $P$ is characteristic. From
$b^2=-2s$ we obtain $\alpha(s)=-s$. From $a=\beta s$ we obtain
$\alpha(\beta)=s^{-1}\beta^{-1}s$ in multiplicative notation.
Follow $\alpha|_P$ by conjugation with $s^{-1}$, then by $\sigma$.
The resulting automorphism $\tau$ fixes $\beta$ and sends $s$ to
$s^{-1}$. It therefore inverts every $\beta$-conjugate of $s$.
These generate the abelian group $M$, so $\tau|_M=-1$.
But $\beta^p=z$ gives both $\tau(z)=z$ and $\tau(z)=-z$,
contradicting $|z|=p>2$.
\end{proof}
The rank-three rotation-group correspondence now constructs a chiral
abstract polyhedron with these generators. That geometric correspondence
is retained explicitly; the preceding argument proves the group,
generation, intersection and failure of reversal themselves.


\section{Class-three collection and small kernels}
In a group of class at most three put $c=[x,y]$, $d=[c,x]$ and
$e=[c,y]$. The elements $d,e$ are central. The commutator identities
$[uv,w]=[u,w]^v[v,w]$ and $[u,vw]=[u,w][u,v]^w$
give by induction
\[
[x^k,y]=c^k d^{\binom k2},\qquad [x,y^k]=c^k e^{\binom k2}
\quad(k\geq0).
\]
Collecting all $x$ factors first, then $y,c,d,e$, similarly gives
\[
(xy)^k=x^ky^k c^{-\binom k2}d^{-\binom k3}
                 e^{-2\binom k3-\binom k2}.
\]
For completeness, if $C_k,D_k,E_k$ are the last three exponents,
moving the final $x,y$ to their positions gives
$C_{k+1}=C_k-k$, $D_{k+1}=D_k+C_k$ and
$E_{k+1}=E_k+C_k-k-\binom k2$, with all three initial values zero.
These recurrences give exactly the displayed formula.
In particular, when $p\geq5$ and the commutator subgroup has
exponent $p$, all correction factors disappear in $(xy)^p$.
Weight-three commutators are multiplicative in each variable:
expand the same two commutator identities; every correction has
weight four and vanishes. Thus simultaneous inversion of the first
two variables inverts each of $[x,y,x]$ and $[x,y,y]$.
Also $[x^{-1},y^{-1}]^{xy}=[x,y]$, so the new weight-two
commutator has the same order.

\begin{lemma}\label{lem:derived}
Let $p$ be an odd prime. If a nonabelian two-generated group $P$ has order $p^3$ or $p^4$,
then $P'$ has exponent $p$ and the nilpotency class is at most three.
If the class is three, then
\[
|P'|=p^2,\quad Z(P)=\gamma_3(P)\cong C_p,\quad P/P'\cong C_p^2.
\]
\end{lemma}
\begin{proof}
A nonabelian finite $p$-group has noncyclic abelianization. For example,
if $P/P'$ were cyclic, its elementary abelian Frattini quotient would
be cyclic and the Burnside basis theorem would make $P$ cyclic.
Thus $|P'|\leq p^2$. Every nontrivial step in the lower central
series strictly decreases order, giving class at most three.
In class two, if $P'$ had order and exponent $p^2$, then
$|Z(P)|\geq p^2$ and $P/Z(P)$ would have order at most $p^2$.
It cannot be cyclic, so is $C_p^2$; hence $x^p\in Z(P)$ for all
$x$, and $[x,y]^p=[x^p,y]=1$. Since a two-generator class-two
commutator group is cyclic, this is a contradiction.

In class three necessarily $|P|=p^4$, $|P'|=p^2$,
$|\gamma_3(P)|=p$, and $P/P'=C_p^2$.
Also $|Z(P)|=p$: a center of order at least $p^2$ would give a
quotient of order at most $p^2$, hence an abelian quotient and class
at most two. Therefore $Z(P)=\gamma_3(P)$.
Suppose $P'$ were cyclic of order $p^2$. For a generating pair $x,y$,
write $c=[x,y]$. Its image generates $P'/\gamma_3(P)$,
so $c$ generates $P'$. Write
\[
c^x=c^{1+\lambda p},\quad c^y=c^{1+\mu p},\quad
x^p=c^a,\quad y^p=c^b.
\]
The conjugation exponents have the displayed shape because the
conjugation action is a $p$-group. Class-three collection and
$p\mid\binom p2$ give $[x^p,y]=c^p$ and $[x,y^p]=c^p$;
all correction terms are powers of $c^{p^2}$.
Comparison with the displayed powers yields
$a\mu\equiv1$ and $-b\lambda\equiv1\pmod p$.
But $[x^p,x]=1$ gives $a\lambda\equiv0\pmod p$, a contradiction.
Thus $P'\cong C_p^2$, proving the assertion.
\end{proof}


\section{The previously proved cyclic lift}
\begin{theorem}[A cyclic lift with an invariant kernel]\label{UF.CyclicLift}
Let $p$ be odd and let $(G_0;a_0,b_0)$ be a chiral rotation pair
with $|a_0|=p^2$, $|b_0|=2p$. Suppose a surjection
$\phi:G_0\to C_p=\gen c$ sends $a_0,b_0$ to $c,c^{-1}$.
For $C_{p^2}=\gen u$ the subgroup generated by
$a=(a_0,u)$ and $b=(b_0,u^{-1})$ in $G_0\times C_{p^2}$
is a chiral rotation group of order $p|G_0|$ and type
$\{p^2,2p^2\}$.
\end{theorem}
\begin{proof}
The two generators lie in the fibre product
$F=\{(g,u^j):\phi(g)=c^j\}$. Projection has kernel
$\gen{(1,u^p)}$, of order $p$. The element
$b^{2p}=(1,u^{-2p})$ generates that kernel, and the marked subgroup
projects onto $G_0$. It is therefore $F$, of order $p|G_0|$.
The direct-product orders give the stated types and $|ab|=2$.
Projection is injective on $\gen a$, since its image has the same
order $p^2$. A common element of $\gen a$ and $\gen b$ maps to
one, hence is one. Finally any simultaneous reversal preserves
the explicitly identified kernel $\gen{b^{2p}}$ and induces a
reversal on $G_0$, a contradiction. This argument does not assume
that arbitrary chiral covers stay chiral, nor that this kernel is
characteristic under every automorphism.
\end{proof}
This is the earlier central-cyclic-lift proof, included in full with
its original scope. Its use below carries no new question or
historical-priority claim.

\section{Attribution and formal scope}
The first three sections retain the accepted Topic F arguments and
attribution: the rank-three test and Klein-four orbit structure
reorganize Hou, Feng, Leemans and Qu, and the integral cyclotomic
arithmetic is classical. The complete split construction, previously
in the first application, is shared here because the second
application uses two of its regular families.
The cyclic-carry theorem reorganizes the constructions of Kong,
Feng, Hou, Leemans and Qu. The small-kernel lemma replaces the
classification and finite small-prime lookup in their tightness
proof; the earlier cyclic lift is included with its exact kernel
hypothesis. No historical priority claim is made for these results.

The classical regular and chiral group-to-polyhedron correspondences,
Burnside's basis theorem and elementary finite-group facts remain
explicit background. The applications separately identify their
retained classification inputs; none is inferred from a finite sample.
Five unchanged declarations in \texttt{OddKernelIntersection} cover
prime-index intersections and reflection identities. Four unchanged
\texttt{LiftCertificate} declarations cover kernel/intersection
transport (three are used by the cyclic lift). Three new
\texttt{CyclicCarryObstruction} declarations prove inversion on an
abelian generated subgroup and the fixed-power/inverted-carry
contradiction. The ring, group constructions, collection identities,
small-order theorem and geometric correspondences remain written,
not formally certified application instances.
\begin{thebibliography}{9}
\bibitem{HFLQ} D.-D. Hou, Y.-Q. Feng, D. Leemans and H.-P. Qu,
\emph{String C-groups of order $4p^m$}, arXiv:2407.10388v1, 2024.
\bibitem{KFH} T.-T. Kong, Y.-Q. Feng, D.-D. Hou, D. Leemans and H.-P. Qu,
\emph{Chiral 3-polytopes of order $2p^m$}, 22-page accepted manuscript,
19 January 2026, DOI10.26493/1855-3974.3716.b64.
\end{thebibliography}\end{document}
